% ==========================================
% SECTION 3: TASK 2 - NAVIGATION
% ==========================================
\section{Task 2: Navigation}
\label{sec: Task 2}

\begin{frame}{Task 2: Autonomous Map Navigation}

    \begin{block}{Objective}
        Perform a \textbf{full autonomous navigation} of a \textbf{known environment} using the previously generated 2D occupancy grid map to find and record the pose of specific Aruco markers.
        \begin{itemize}
            \item \textbf{Pick Station Marker:} Aruco Marker ID 26, size 0.25m.
            \item \textbf{Place Station Marker:} Aruco Marker ID 238, size 0.25m.
        \end{itemize}
    \end{block}

    \vspace{0.1cm}

    \begin{columns}
        \begin{column}{0.45\textwidth}
            \begin{figure}
                \centering
                % Screenshot di Gazebo con il robot che naviga verso un marker ArUco
                \includegraphics[width=0.7\linewidth]{PickTable.png}
                \caption{Pick Station Marker 26.}
            \end{figure}
        \end{column}

        \begin{column}{0.45\textwidth}
            \begin{figure}
                \centering
                % Screenshot di Gazebo con il robot che naviga verso un marker ArUco
                \includegraphics[width=0.7\linewidth]{PlaceTable.png}
                \caption{Place Station Marker 238.}
            \end{figure}
        \end{column}
    \end{columns}

\end{frame}

\begin{frame}{Task 2: Navigation Trade-offs}

    \begin{block}{Operational Challenges}
        ArUco markers are attached to tables that may be located in corners or occluded areas.
    \end{block}

    \begin{center}
        \large
        \textbf{Strategic Trade-off: Search Logic}
    \end{center}

    \begin{alertblock}{Exhaustive Perimeter Search}
        To locate the tables, we evaluated an \textit{exhaustive perimeter search} (wall-following) to systematically scan every corner.
    \end{alertblock}
    
    \vspace{0.1cm}

    \begin{exampleblock}{Probabilistic Heuristic (Naive Explorer)}
        However, we opted for a \textbf{probabilistic heuristic}. By generating safe random goals within the free space, the robot covers the environment more rapidly than tracing the full perimeter, maintaining high detection reliability with lower execution time.
    \end{exampleblock}

\end{frame}

% Slide 1: System Orchestration
\begin{frame}{Task 2: System Architecture \& Orchestration}

    The complex behavior required for Task 2 is managed by a central \textbf{Coordinator Node} that acts as the high-level decision maker.

    \begin{columns}
        \begin{column}{0.5\textwidth}
            \begin{block}{The Coordinator Design}
                \begin{itemize}
                    \item \textbf{Finite State Machine (FSM):} Strict transition logic between \textit{Exploring}, \textit{Tracking}, and \textit{Inspecting} states.
                    \item \textbf{Concurrency Handling:} Utilizes \texttt{ReentrantCallbackGroup} and \texttt{MultiThreadedExecutor}. This allows the node to process sensory interrupts (e.g., seeing a marker) immediately, even while blocking navigation actions are active.
                    \item \textbf{Interrupt Logic:} The standard exploration loop is designed to be preempted instantly upon ArUco detection.
                \end{itemize}
            \end{block}
        \end{column}

        \begin{column}{0.4\textwidth}
            \begin{figure}
                \centering
                % Qui starebbe benissimo un diagramma a stati (Init -> Exploring <-> Tracking -> Homing)
                \includegraphics[width=\linewidth]{task2_state_machine.png}
                \caption{State Machine Logic}
            \end{figure}
        \end{column}
    \end{columns}

\end{frame}

% Slide 2: Global Localization
\begin{frame}{Task 2: Global Localization}
    Since the robot starts in an unknown location on the map, we implemented a global localization routine exploiting the Adaptive Monte Carlo Localization (AMCL).
    \begin{columns}[c]
        \begin{column}{0.6\textwidth}
            \begin{block}{The Converge Strategy:}
                \begin{enumerate}
                    \item \textbf{Particle Spread:} Trigger \texttt{/reinitialize\_global\\\_localization} to distribute AMCL particles uniformly across the map.
                    \item \textbf{Data Gathering:} The robot executes a $360^\circ$ in-place rotation at $1.0$ rad/s to feed the laser scan matcher.
                    \item \textbf{Covariance Check:} The routine continuously monitors the pose covariance matrix, ensuring it falls below strict thresholds before allowing any movement.
                \end{enumerate}
            \end{block}
        \end{column}
        
        \begin{column}{0.3\textwidth}
            \begin{figure}
                \centering
                % Screenshot di RViz con le particelle AMCL sparse
                \includegraphics[width=0.7\linewidth]{Convergence2.png}
            \end{figure}

            \begin{figure}
                \centering
                % Grafico che mostra la diminuzione della covarianza nel tempo durante l'inizializzazione
                \includegraphics[width=1.0\linewidth]{covariance_plot.png}
                \caption{Covariance Reduction Over Time}
            \end{figure}
        \end{column}
    \end{columns}

\end{frame}

% Slide 3: Exploration Strategy
\begin{frame}{Task 2: Exploration Strategy}

    After the localization, the Coordinator Node activate the following nodes:
    \begin{columns}
        \begin{column}{0.5\textwidth}
            \begin{block}{Generalized Nodes}
                \begin{itemize}
                    \item \textbf{Aruco Marker Detector:} The node continuously monitors the camera feed for the station markers.
                    \item \textbf{Explorator:} The node processes the occupancy grid map to create a "whitelist" of safe navigation targets, ensuring the robot only attempts to navigate to locations that are free of obstacles.
                    \item \textbf{Navigator:} The robot continuously samples from the "whitelist" of safe exploration goals. After reaching each goal, the robot performs a $360^\circ$ rotation to improve its environmental awareness.
                    \end{itemize}
            \end{block}
        \end{column}
            
        \begin{column}{0.5\textwidth}
            \begin{figure}
                \centering
                \includegraphics[width=0.8\linewidth, height=2.5cm, keepaspectratio]{exploration_map.png}
                \caption{Clearance-Based Goal Sampling}
            \end{figure}

            \begin{figure}
                \centering
                \includegraphics[width=0.8\linewidth, height=2cm, keepaspectratio]{RandomNav2.png}
                \caption{Random Safe Goal Navigation}
            \end{figure}
        \end{column}
    \end{columns}
    
    \vspace{0.1cm}

\end{frame}


\begin{frame}{Task 2: Exploration Strategy - Contd.}
    
    \begin{block}{Exploration Strategy}
        \begin{itemize}
            \item \textbf{Tracking:} Once an ArUco marker is detected, the robot is brought to a specific distance from the marker, while continuously facing it by adjusting the camera tilt and pan angles.
            \item \textbf{Inspecting:} Once the robot is correctly positioned, it transitions to the Inspecting state, where the marker and robot Pose is saved in a csv file for later use.
            \item \textbf{Homing:} After all the markers are detected, the robot automatically navigates back to the Home Position.
        \end{itemize}
    \end{block}

    \begin{columns}
        
        \begin{column}{0.45\textwidth}
            \begin{figure}
                \centering
                % Screenshot di RViz con il robot che segue un marker ArUco con la  testa bassa e la camera che inquadrano il marker
                \includegraphics[width=0.8\linewidth, height=3cm, keepaspectratio]{ArucoTracking2.png}
                \caption{Aruco Tracking Behavior}
            \end{figure}
        \end{column}

        \begin{column}{0.45\textwidth}
            \begin{figure}
                \centering
                % Screenshot di RViz con il robot che guarda il marker arrivato alla stazione
                \includegraphics[width=0.8\linewidth, height=3cm, keepaspectratio]{ArucoInspection.png}
                \caption{Aruco Inspecting Behavior}
            \end{figure}
        \end{column}
        
    \end{columns}

\end{frame}

% Forse questa slide si può far volare o metterla nell'appendice
\begin{frame}{Task 2: Safety Navigation}
    \begin{block}{Map Processing Pipeline}
            Using OpenCV on the occupancy grid map:
            \begin{enumerate}
                \item \textbf{Safety Mapping:} Applied \texttt{cv2.distanceTransform} to calculate the distance of every pixel to the nearest obstacle.
                \item \textbf{Filtering:} Created a "Whitelist Mask" of pixels where:
                $$ d_{obstacle} > (R_{robot} + \text{Clearance}_{0.7m}) $$
                \item \textbf{Goal Selection:} Randomly samples coordinates from the whitelist and converts them from Image Frame (pixels) to World Frame (meters).
            \end{enumerate}
    \end{block}

        \begin{center}
            \textit{Result: The robot never attempts to navigate into walls or narrow traps.}
        \end{center}
\end{frame}

\begin{frame}{Task 2: Navigation Video Result}
    \begin{figure}
        \centering
        % Ho aggiunto \ davanti a tutti i % nel link qui sotto:
        \href{https://liveunibo.sharepoint.com/sites/idughis/_layouts/15/stream.aspx?id=\%2Fsites\%2Fidughis\%2FShared\%20Documents\%2FGeneral\%2FMagistrale\%2FAMR\%2FVideos\%2FTask2\%2Emp4&referrer=StreamWebApp\%2EWeb&referrerScenario=AddressBarCopied\%2Eview\%2E3889980f\%2D8794\%2D46a6\%2Db554\%2D9371fc06f370}{
            % Ho cambiato 0.2 in 0.7, altrimenti l'immagine è minuscola
            \includegraphics[width=0.9\linewidth]{Thumb_task2.png}
        }
        \caption{Navigation Task Execution (Click to play in browser)}
    \end{figure}
\end{frame}